\paragraph{Graduación, unidad y automorfismos.}
\label{inc-ampl-automorfismos}
La construcción de Frenkel--Lepowsky--Meurman se aplica al retículo
$N_\alpha$ una vez verificadas sus propiedades. Incorpora el espacio
de osciladores, la extensión central del retículo, el producto de vértice
y el sector torcido por el levantamiento de la negación. El resultado es
el álgebra $V^\natural_{\rm HMT}$ de \eqref{exc:flm}, cuyo grupo
de automorfismos es isomorfo al Monstruo. Su dominio de acción queda así
fijado por una estructura algebraica con graduación y producto. El teorema
de Moonshine de Borcherds determina posteriormente las trazas graduadas
con inserción de esos automorfismos \cite{flm,borcherds}.

La unidad adquiere en este nivel una realización precisa: la componente
de peso cero es la línea del vacío y aporta el coeficiente $1$ de
$q^{-1}$ en el carácter graduado. En el nivel cilíndrico, las
aplicaciones de refinamiento conservan la función constante mediante
precomposición. Cada afirmación de conservación tiene, por tanto, su
objeto y su aplicación. El desarrollo de las pruebas hace visible esa
continuidad de construcción, desde las operaciones aritméticas y la
memoria hasta la incidencia, la métrica y los productos de vértice.

Esta articulación sitúa a $K$ y a la clausura de $\alpha$ en el argumento
principal. $K$ conserva reversiblemente una coordenada integral de la
historia; la normalización de $\alpha$ compone los registros regionales
con esa coordenada; la configuración firmada determina una bandera; y
la bandera marcada selecciona la modificación reticular sobre la que se
realiza Moonshine. Las demostraciones siguientes desarrollan estas
aplicaciones, distinguiendo las equivalencias reversibles, los cocientes
de información y las realizaciones que añaden métrica o estructura
algebraica. Así se expresa una relación demostrable entre estructura,
ecuación y valor, conservando la atribución de cada construcción y el
alcance de sus teoremas.
